\begin{tikzpicture}[
    >=Latex,
    scale=1.4,
    every node/.style={font=\small},
]
    % =========================================================
    % T-Trager mit vertikalem viridis-aehnlichem Gradient
    % Stamm: x in [-0.3, 0.3], y in [0, 2.5]
    % Querbalken: x in [-2, 2], y in [2.5, 3.0]
    % =========================================================
    \begin{scope}
        \clip
            (-0.3, 0) -- (0.3, 0) -- (0.3, 2.5) --
            (2, 2.5) -- (2, 3.0) --
            (-2, 3.0) -- (-2, 2.5) --
            (-0.3, 2.5) -- cycle;
        \shade[bottom color=viridisLow, top color=viridisHigh,
            middle color=viridisMid]
            (-2.1, -0.05) rectangle (2.1, 3.05);
    \end{scope}

    % --- T-Trager Aussenkontur ---
    \draw[thick, black]
        (-0.3, 0) -- (0.3, 0) -- (0.3, 2.5) --
        (2, 2.5) -- (2, 3.0) --
        (-2, 3.0) -- (-2, 2.5) --
        (-0.3, 2.5) -- cycle;

    % =========================================================
    % Iso-Linien (Phi = const)
    % =========================================================
    % Im Stamm: drei horizontale Linien (Phi = 0.2, 0.4, 0.6)
    \draw[white, very thick] (-0.3, 0.6) -- (0.3, 0.6);
    \node[white, font=\scriptsize, anchor=west] at (0.4, 0.6)
        {$\Phi=0{,}2$};

    \draw[white, very thick] (-0.3, 1.2) -- (0.3, 1.2);
    \node[white, font=\scriptsize, anchor=west] at (0.4, 1.2)
        {$\Phi=0{,}4$};

    \draw[white, very thick] (-0.3, 1.8) -- (0.3, 1.8);
    \node[white, font=\scriptsize, anchor=west] at (0.4, 1.8)
        {$\Phi=0{,}6$};

    % Im Querbalken: gebogene Iso-Linie (Phi = 0.85), zeigt das
    % Aufweiten in die Querbalkenarme
    \draw[white, very thick]
        (-1.8, 2.55) .. controls (-0.6, 2.95) and (0.6, 2.95) .. (1.8, 2.55);
    \node[white, font=\scriptsize, anchor=north east] at (1.85, 2.55)
        {$\Phi=0{,}85$};

    % =========================================================
    % Gradient-Pfeile (+nabla Phi, zeigt nach hoher Phi = oben)
    % =========================================================
    \foreach \x/\y in {0/0.9, 0/2.1, 1.4/2.75, -1.4/2.75} {
        \draw[->, very thick, white] (\x, \y) -- (\x, \y+0.25);
    }
    \node[white, font=\scriptsize, anchor=west] at (0.05, 0.95)
        {$+\nabla\Phi$};

    % =========================================================
    % Randbeschriftungen
    % =========================================================
    \draw[<-, thin] (0, -0.05) -- (0, -0.5);
    \node[anchor=north, font=\scriptsize, align=center] at (0, -0.5)
        {$\Gamma_\text{unten}$\\$\Phi = 0$};

    \draw[<-, thin] (0, 3.05) -- (0, 3.5);
    \node[anchor=south, font=\scriptsize, align=center] at (0, 3.5)
        {$\Gamma_\text{oben}$\\$\Phi = 1$};

    % =========================================================
    % Farbskala (rechts neben dem Bauteil)
    % =========================================================
    \begin{scope}[xshift=3.4cm]
        \shade[bottom color=viridisLow, top color=viridisHigh,
            middle color=viridisMid]
            (0, 0) rectangle (0.3, 3.0);
        \draw[thick] (0, 0) rectangle (0.3, 3.0);

        \foreach \v/\l in {0/0, 0.75/0{,}25, 1.5/0{,}5,
                        2.25/0{,}75, 3.0/1} {
            \draw (0.3, \v) -- (0.4, \v);
            \node[font=\scriptsize, anchor=west] at (0.45, \v) {\l};
        }
        \node[font=\scriptsize, anchor=south] at (0.15, 3.05) {$\Phi$};
    \end{scope}

\end{tikzpicture}